\section{Evaluation Results}

\subsection{Evaluation Overview}

This section reports the executed proof-of-concept evaluation for the dataset \texttt{openaq\_capitals\_2025\_h2}. The results should be read as a controlled threat-coverage evaluation of integrity mechanisms, not as an environmental analysis of the monitored locations. The experiment compares four integrity models against a fixed set of tampering scenarios and records whether the verifier produced the expected alert code for each applicable model/threat pair.

Table~\ref{tab:dataset-summary} summarizes the dataset preparation used for the evaluation. The preprocessing stage converted raw OpenAQ records into canonical measurements and excluded records with missing measurement values.

\begin{table}[htbp]
\centering
\caption{Dataset preparation summary for \texttt{openaq\_capitals\_2025\_h2}.}
\label{tab:dataset-summary}
\begin{tabular}{lr}
\toprule
Item & Value \\
\midrule
Input OpenAQ records & 116,361 \\
Canonical records after preprocessing & 112,973 \\
Dropped records & 3,388 \\
Selected monitoring locations & 12 \\
Cities & 4 \\
Monitoring locations per city & 3 \\
Parameters after preprocessing & 9 \\
\bottomrule
\end{tabular}
\end{table}

\subsection{Compared Integrity Models}

Table~\ref{tab:model-summary} summarizes the four compared integrity models. The model sequence is designed to isolate the effect of adding an audit trail, then hash-chain continuity, and finally provenance and permission reconstruction.

\begin{table}[htbp]
\centering
\caption{Integrity models compared in the proof-of-concept.}
\label{tab:model-summary}
\small
\begin{tabular}{lp{0.62\linewidth}}
\toprule
Model & Verification capability \\
\midrule
A & Conventional storage with schema and duplicate record identifier checks. \\
B & Audit trail with event payload and event identifier checks, without hash-chain continuity. \\
C & Audit trail plus hash chain, adding previous-hash and block-hash continuity checks. \\
D & Hash chain plus provenance and permissions, adding actor/key, permission, correction, provenance, and delayed-synchronization checks. \\
\bottomrule
\end{tabular}
\end{table}

\subsection{Executed Scenario Matrix}

The full implemented scenario matrix generated 25 controlled tampering scenarios. For each scenario, the proof-of-concept produced a tampered artifact, a scenario label, a verification output, and a per-scenario evaluation output. The aggregate status counts are shown in Table~\ref{tab:status-counts}.

The run produced 20 detected cases and 5 expected non-detections. No scenarios were marked as missed, partial, or unexpected-alert cases. These counts should be interpreted as scenario-level verification outcomes for the controlled benchmark, not as statistical detection rates.

\begin{table}[htbp]
\centering
\caption{Aggregate scenario status counts.}
\label{tab:status-counts}
\begin{tabular}{lr}
\toprule
Status & Count \\
\midrule
\texttt{detected} & 20 \\
\texttt{expected\_not\_detected} & 5 \\
\texttt{missed} & 0 \\
\texttt{partial} & 0 \\
\texttt{unexpected\_alert} & 0 \\
\bottomrule
\end{tabular}
\end{table}

\subsection{Threat-Coverage Matrix}

Table~\ref{tab:threat-coverage} presents the threat-coverage matrix for Models A-D. The matrix uses three compact status codes: D for detected, END for expected non-detection, and NA for not applicable. A detected cell indicates that the expected alert code was present in the verifier output. An expected non-detection indicates an explicit model limitation rather than an implementation failure. A not-applicable cell indicates that the scenario was outside the implemented scope of the model.

\begin{table}[htbp]
\centering
\caption{Threat-coverage matrix for integrity Models A-D.}
\label{tab:threat-coverage}
\small
\begin{tabular}{p{0.42\linewidth}cccc}
\toprule
Threat & A & B & C & D \\
\midrule
\texttt{value\_modification} & END & D & D & D \\
\texttt{timestamp\_modification} & END & D & D & D \\
\texttt{record\_deletion} & END & END & D & D \\
\texttt{fake\_record\_insertion} & END & D & D & D \\
\texttt{replay} & D & D & D & D \\
\texttt{broken\_provenance} & NA & NA & NA & D \\
\texttt{unauthorized\_correction} & NA & NA & NA & D \\
\texttt{revoked\_actor\_key\_usage} & NA & NA & NA & D \\
\texttt{missing\_correction\_reason} & NA & NA & NA & D \\
\texttt{delayed\_synchronization} & NA & NA & NA & D \\
\bottomrule
\end{tabular}
\end{table}

The model-level coverage counts are summarized in Table~\ref{tab:applicable-coverage}. Not-applicable scenarios are excluded from the applicable-scenario denominator. Expected non-detections are retained as explicit model limitations.

\begin{table}[htbp]
\centering
\caption{Applicable scenario coverage by model.}
\label{tab:applicable-coverage}
\begin{tabular}{lrrrr}
\toprule
Model & Applicable & Detected & Expected not detected & Other \\
\midrule
A & 5 & 1 & 4 & 0 \\
B & 5 & 4 & 1 & 0 \\
C & 5 & 5 & 0 & 0 \\
D & 10 & 10 & 0 & 0 \\
\bottomrule
\end{tabular}
\end{table}

\subsection{Model-Level Interpretation}

Model A detected only the replay scenario in the implemented matrix because replay introduced a duplicate record identifier. It did not contain independent audit or hash evidence for detecting direct value modification, timestamp modification, record deletion, or fake insertion.

Model B detected value modification, timestamp modification, fake record insertion, and replay through audit-event integrity checks. It did not detect record deletion in this implementation because its audit events were not linked into a hash chain.

Model C detected all five base scenarios. The addition of hash-chain linkage made stream-continuity disruption visible, including record deletion.

Model D detected all Model C scenarios and the five additional provenance, permission, correction, revoked-key, and delayed-synchronization scenarios. This result supports the intended role of Model D as the most expressive model in the proof-of-concept, while remaining limited to the controlled scenarios implemented in this benchmark.

\subsection{Reproducibility Artifacts}

The primary machine-readable outputs are the scenario metrics CSV file, the threat-coverage matrix CSV file, the metrics summary JSON file, and the experiment-run manifest. The manifest records artifact paths, file sizes, hashes, scenario counts, and run metadata for the executed scenario matrix.
